%\documentclass[]{article}
%\usepackage[spanish]{babel}
%\usepackage[utf8]{inputenc}
%\usepackage{geometry}
%\usepackage{colortbl}
%\usepackage{longtable}
%\usepackage{graphicx}
%\geometry{tmargin=3cm,bmargin=3cm,lmargin=3cm,rmargin=2cm}
%
%\begin{document}
\begin{center}
\begin{longtable}{|p{0.225\textwidth}|p{0.225\textwidth}|p{0.225\textwidth}|p{0.225\textwidth}|}
\hline
{\bf {Empresa:}} &
\multicolumn{2}{p{0.45\textwidth}|} { Escuela de Ingeniería de Sistemas y Computación } &
{\includegraphics[width=80.5pt]{LOGO}} \\
\hline
\bf {Nombre del caso de uso:}&\multicolumn{3}{l|}{
Consulta Avanzada.
} \\
\hline
\bf Codigo: & 
CU12 &\bf Fecha: & 
Abril 02 2011 \\
\hline
\bf Autor(es ): & 
Edgar Andrés Moncada & 
 & 
 \\
\hline
\bf Descripcion: &\multicolumn{3}{p{0.675\textwidth}|}
{
Permite la realización de consulta de los documentos digitales del Sistema de manera Avanzada.
} \\
\hline
\bf Actores: &\multicolumn{3}{p{0.675\textwidth}|}{
Administrador, Catalogador, Usuario Normal, Usuario No Registrado 
} \\
\hline
\bf Precondiciones: &\multicolumn{3}{p{0.675\textwidth}|}
{
El Sistema se haya iniciado.
} \\
\hline
\multicolumn{4}{|c|}{\bf {Flujo Normal}}\\
\hline
\multicolumn{2}{|c}{\bf Actor} & \multicolumn{2}{|c|}{\bf Sistema } \\
\hline
\multicolumn{2}{|p{0.45\textwidth}}
{
\begin{itemize}
\item[1. ]El caso de uso inicia cuando el Usuario va a realizar una consulta y selecciona la opción de Consulta Avanzada.
\item[3.]El Usuario ingresa datos en alguno o en combinaciones de los campos de texto, seleccionando los parámetros deseados y selecciona la opción consultar.
\item[6. ]El Usuario selecciona uno de los resultados de la consulta. 
\end{itemize}
} &
\multicolumn{2}{|p{0.45\textwidth}|}
{
\begin{itemize}
\item[2.]El Sistema muestra una interface 3 campos de texto para Titulo, Autor y Palabra Clave con la posibilidad de seleccionar para cada una las restricciones de con todas las palabras, con algunas de las palabras o con ninguna de las palabras; Incluye también los parámetros de Área, Idioma, Fecha de Publicación y Formato del Archivo.
\item[4. ]El Sistema realiza la consulta.
\item[5. ]El Sistema muestra un listado de los resultados de la consulta.
\item[7. ] El Sistema muestra todos los datos relacionados con el documento y el caso de uso termina.
\end{itemize}
} \\
\hline
\multicolumn{4}{|c|}{\bf {Flujo Alternativo}}\\
\hline
\multicolumn{2}{|p{0.45\textwidth}}
{
\begin{itemize}
 \item[6.2. ] El Usuario no selecciona ninguno de los resultados de la consulta y el caso de uso termina.
\end{itemize}
} &
\multicolumn{2}{|p{0.45\textwidth}|}
{
\begin{itemize}
 \item[5.1. ] La consulta no arroja ningún resultado.
 \item[6.1. ] El Sistema envía una notificación indicando que no se encontraron coincidencias.
 \end{itemize}
} \\
\hline
\bf Poscondiciones: &\multicolumn{3}{p{0.675\textwidth}|}
{
Se muestra las consultas relacionadas con los datos que ingreso el Usuario.
} \\
\hline
\bf Excepciones: &\multicolumn{3}{p{0.675\textwidth}|}
{
No se pude acceder a la base de datos para obtener los datos de los Documentos Digitales.
} \\
\hline
\bf Aprobado por : & 
 & \bf Fecha & 
 \\
\hline
\end{longtable}
\end{center}
%\end{document}